% 第1章 绪论
\chapter{绪论}
\label{chap:intro}

\section{研究背景与意义}
\label{sec:bg}

光与物质的相互作用是现代物理学的核心议题之一。在纳米尺度上，电磁场在单个光学周期内经历复杂的时间演化过程，其特征时间尺度为飞秒（$\SI{1}{fs} = 10^{-15}\;\text{s}$）量级，特征空间尺度为纳米至亚纳米量级。理解并直接测量这一时空尺度上的电磁场分布$\mathbf{E}(\mathbf{r}, t)$，对于纳米光子学、超快光谱学和量子光学等领域具有重要的科学意义。

传统的光学测量方法受到衍射极限的制约，空间分辨率被限制在半波长量级（$\sim\lambda/2$）。以可见光波段为例，800\,nm波长的光其衍射极限约为400\,nm，远远无法满足纳米尺度空间分辨的需求。近场光学技术（如散射型扫描近场光学显微镜，s-SNOM）可将空间分辨率提升至$\sim\SI{10}{nm}$，但其时间分辨率受限于探测脉冲的持续时间，通常为数十飞秒，不足以分辨单个光场周期内的动力学过程。

电子显微镜在空间分辨率方面具有天然优势。keV能量的电子具有皮米（$\SI{1}{pm} = 10^{-12}\;\text{m}$）量级的德布罗意波长，远超光学衍射极限。然而，传统超快电子显微术的时间分辨率受限于电子脉冲宽度，由于库仑排斥和空间电荷效应，脉冲宽度停留在$\sim\SI{100}{fs}$，比可见光的光学周期（$\sim\SI{2.7}{fs}$ @ 800\,nm）长两个数量级。这一"时间分辨率的鸿沟"阻碍了对亚周期电磁动力学的直接观测。

阿秒科学的兴起为突破这一时间壁垒提供了全新的工具。2001年，Hentschel等人\cite{Hentschel2001}和Drescher等人\cite{Drescher2001}首次在实验上产生了孤立阿秒脉冲，打开了亚飞秒时间分辨的大门。2023年诺贝尔物理学奖授予Agostini、Krausz和L'Huillier三位科学家，以表彰他们在阿秒脉冲产生方面的开创性贡献\cite{nobel2023}。然而，光子基的阿秒探针受制于光波长的衍射极限，即使在30\,nm的软X射线波段，空间分辨率也仅能达到$\sim\SI{15}{nm}$——比keV电子的皮米级分辨率粗糙三个数量级。

本文所研究的阿秒电子显微术，将阿秒科学的时间分辨能力与电子显微镜的空间分辨能力有机结合，填补了超快科学中高时空分辨测量的关键空白。图\ref{fig:ch1-resolution-map}展示了不同超快测量技术在时间-空间分辨率图上的定位。

\begin{figure}[htbp]
    \centering
    \includegraphics[width=0.85\linewidth]{ch1/fig_resolution_map}
    \caption{\label{fig:ch1-resolution-map} 超快测量技术的时间-空间分辨率对比。横轴为时间分辨率，纵轴为空间分辨率。阿秒电子显微术（红色星号）位于时间分辨率$\sim\SI{800}{as}$、空间分辨率$\sim\SI{50}{nm}$的独特位置，填补了光子探针和电子探针之间的空白。}
\end{figure}

\section{超快电子显微术的发展}
\label{sec:uem-dev}

超快电子显微术的发展可追溯至Zewail课题组在2000年代初的开创性工作。Zewail等人\cite{Zewail2010}建立了四维电子显微术（4D electron microscopy），利用飞秒激光触发泵浦-探测方案，在透射电子显微镜中实现了$\sim\SI{1}{nm}$空间分辨率和$\sim\SI{100}{fs}$时间分辨率的联用。这一技术使得对结构动力学（如相变、晶格振动、化学反应中间态）的实时观测成为可能。

然而，泵浦-探测型超快电子显微术面临两个根本性限制。第一，电子脉冲的时间展宽：由于电子之间的库仑排斥效应（空间电荷效应），飞秒电子脉冲在传播过程中不可避免地展宽，即使在单电子模式下，脉冲宽度也受到初始能量弥散的限制。第二，时间分辨率的不对称性：时间分辨率受限于泵浦脉冲和探测脉冲中较长的一个，且无法突破光学周期的限制。

针对上述限制，研究者发展了多种电子脉冲压缩技术。Baum和Zewail\cite{Baum2007}提出了利用射频腔对电子脉冲进行时间压缩的方案。Ryabov和Baum\cite{Ryabov2016}利用太赫兹脉冲驱动的电子脉冲压缩，实现了$\sim\SI{1}{fs}$的时间分辨率，并实现了电磁波波形的直接成像。

近年来，超快透射电子显微镜（UTEM）技术也取得了重要进展。Feist等人\cite{Feist2015}在UTEM中实现了光控电子量子态的相干调制，证明了自由电子与光场之间的量子相干相互作用。这一工作为光场驱动的电子波函数操控奠定了实验基础。

\section{阿秒科学的产生与进展}
\label{sec:attosecond-dev}

阿秒科学的根基在于高次谐波产生（high-harmonic generation, HHG）现象。1993年，Corkum\cite{Corkum1993}提出了著名的三步模型（three-step model），将强场电离过程分解为隧道电离、自由电子在激光场中的加速运动、以及电子与母离子的复合辐射三个阶段。这一模型定量解释了HHG的平台区和截止频率，为阿秒脉冲的产生提供了理论指导。

Antoine等人\cite{Antoine1996}从理论上预言了利用HHG产生阿秒脉冲串的可能性。2001年，两个研究组分别在实验上实现了阿秒脉冲的产生和表征：Hentschel等人\cite{Hentschel2001}产生了$\sim\SI{650}{as}$的脉冲，Drescher等人\cite{Drescher2001}实现了$\sim\SI{2.7}{fs}$的软X射线脉冲。这些突破性进展标志着阿秒科学的正式诞生。

Krausz和Ivanov\cite{Krausz2009}在2009年的综述中系统总结了阿秒物理的发展，指出了阿秒脉冲在原子分子动力学观测中的广泛应用前景。随后，阿秒脉冲的持续时间被不断缩短，目前已达到$\sim\SI{43}{as}$的最短记录\cite{Li2020_as}。在空间分辨率方面，Popmintchev等人\cite{Popmintchev2012}将HHG扩展到keV硬X射线波段，为高能量阿秒光子提供了可能。

阿秒科学的应用范围也在不断扩大。从最初的原子物理（电子隧穿观测\cite{Uiberacker2007}），到分子动力学（电子定位\cite{Sansone2010}），再到凝聚态物理（超快退磁\cite{Kirilyuk2010}），阿秒技术正在向越来越多的学科渗透。Varillas等人\cite{Varillas2026}在2026年的阿秒科学路线图中展望了未来十年的发展方向，强调了阿秒技术与电子显微术结合的重要潜力。

\section{阿秒电子显微术的提出与科学问题}
\label{sec:proposal}

阿秒电子显微术的核心理念是将阿秒级时间结构"刻印"到电子束中，使电子束自身成为阿秒时间分辨的探针。这一思想的理论基础源于Tsarev、Ryabov和Baum\cite{Tsarev2021}提出的时间Talbot效应：连续波激光对电子波函数进行周期性能量调制后，电子在自由空间传播过程中会经历Talbot回复，在特定传播距离处形成阿秒级时间宽度的电子脉冲串。

与传统泵浦-探测方案不同，这一方法不需要飞秒激光器：一束连续波（continuous-wave, CW）激光即可产生阿秒电子脉冲。这一特点大大降低了实验复杂度，同时从根本上避免了空间电荷效应导致的脉冲展宽——因为每个光学周期内最多只有一个电子通过样品。

本文聚焦于以下科学问题：

\begin{enumerate}
    \item 如何利用连续波激光在常规透射电子显微镜中产生可调控的阿秒电子脉冲？
    \item 阿秒电子脉冲如何实现对纳米结构近场的亚周期时间分辨成像？
    \item 手性纳米结构中的表面波动力学在亚周期尺度上呈现怎样的特征？
    \item 介质纳米共振器中不同阶 Mie 模式之间的时间延迟是否可以直接观测？
    \item 纳米加工精度如何影响多天线系统的电磁响应对称性？
\end{enumerate}

图\ref{fig:ch1-framework}展示了本文的研究框架和各章节之间的逻辑关系。

\begin{figure}[htbp]
    \centering
    \includegraphics[width=0.9\linewidth]{ch1/fig_framework}
    \caption{\label{fig:ch1-framework} 本文研究框架。从理论基础出发，发展阿秒电子脉冲产生方法，进而应用于三类纳米光子系统（针尖、共振器、超原子）的亚周期动力学研究。}
\end{figure}

\section{本文主要研究内容与章节安排}
\label{sec:outline}

本文共分为七章，各章内容安排如下：

第一章为绪论，介绍研究背景、超快电子显微术和阿秒科学的发展历程，提出本文的科学问题。

第二章为理论基础，系统阐述电子波函数的光场调制机制、时间Talbot效应与阿秒脉冲形成的物理图像、高次谐波产生的三步模型、以及电子能量损失谱在近场成像中的应用原理。

第三章介绍阿秒电子脉冲的产生与表征方法，包括连续波激光驱动的电子脉冲压缩方案、能量滤波器的对比度增强机制、时间扫描与延迟控制的实现方式、以及时间分辨率的标定方法。

第四章至第六章分别展示三类纳米光子系统的实验研究结果：针尖纳米结构的手性表面波动力学、介质纳米共振器的多极模式干涉、以及多天线超原子的对称性破缺响应。

第七章总结全文工作，归纳主要创新点，并对阿秒电子显微术的未来发展方向进行展望。
